\bookchapter{Git: Branches, Merging and Conflicts}\label{ch:git2}

\section*{Explanation}
You want to try a faster walking gait, but the robot must keep walking with the working code until the new idea works. The solution is a \textbf{branch}: a parallel line of development. You experiment on the branch, and if it works you \textbf{merge} it back. If not, you delete it and nothing is lost. Branches in Git are cheap (a branch is only a label, see Figure~\ref{fig:g2_branch}), so professionals use one for every idea, bug fix and feature.

\section*{Branches are labels}
Recall from Chapter~\ref{ch:git1}: \texttt{main} is a label on a commit. A new branch is a second label. \texttt{HEAD} points to the branch that you are on, and a new commit moves \emph{that} label forward.
\begin{center}
\begin{tikzpicture}[c/.style={draw,circle,thick,minimum size=1cm,fill=green!12,font=\scriptsize\ttfamily},lb/.style={draw,rounded corners,font=\footnotesize\ttfamily},>={Stealth[length=2.5mm]}]
  \node[c] (a) at (0,0) {A};
  \node[c] (b) at (1.8,0) {B};
  \node[c] (cc) at (3.6,0) {C};
  \node[c,fill=orange!25] (d) at (5.4,0.9) {D};
  \node[c,fill=orange!25] (e) at (7.2,0.9) {E};
  \node[c] (f) at (5.4,-0.9) {F};
  \draw[->,thick] (b)--(a); \draw[->,thick] (cc)--(b); \draw[->,thick] (d)--(cc); \draw[->,thick] (e)--(d); \draw[->,thick] (f)--(cc);
  \node[lb,fill=yellow!30] (m) at (5.4,-2.0) {main};
  \node[lb,fill=yellow!30] (ft) at (8.9,0.9) {feature/speed};
  \draw[->,thick] (m)--(f); \draw[->,thick] (ft)--(e);
  \node[lb,fill=blue!10] (h) at (3.4,-2.0) {HEAD};
  \draw[->,thick,dashed] (h)--(m);
\end{tikzpicture}
\captionof{figure}{Two branches that have diverged after commit C. \texttt{feature/speed} has the commits D and E, \texttt{main} got commit F in the meantime. HEAD is on \texttt{main}.}\label{fig:g2_branch}
\end{center}

\section*{Working with branches}
\begin{verbatim}
$ git branch                       # list branches, * = current
* main
$ git switch -c feature/speed      # create and switch to a new branch
Switched to a new branch 'feature/speed'
$ git switch main                  # go back
$ git branch -vv                   # with last commit and upstream
$ git branch -d feature/speed      # delete (only if merged)
$ git branch -D feature/speed      # delete even if NOT merged
\end{verbatim}
(Older tutorials use \texttt{git checkout}; \texttt{git switch} for branches and \texttt{git restore} for files are the clearer modern replacements.) When you switch, Git rewrites the files in your working directory to match that branch. Therefore switch only when \texttt{git status} is clean or the changes are stashed (below).

Use clear names: \texttt{feature/gait-speed}, \texttt{fix/camera-timeout}, \texttt{docs/readme}. The slash only groups the names.

\section*{Merging}
\begin{verbatim}
$ git switch -c feature/speed
$ # ... edit, add, commit ...
$ git switch main
$ git merge feature/speed
\end{verbatim}
\texttt{git merge X} brings the commits of branch \texttt{X} into the branch you are on. There are two cases (Figure~\ref{fig:g2_merge}).
\begin{description}
  \item[Fast-forward.] If \texttt{main} has not moved since you branched, Git just moves the label \texttt{main} forward; the history stays a straight line.
  \item[Merge commit.] If both branches have new commits, Git combines the changes and creates a commit with \emph{two} parents.
\end{description}
\begin{verbatim}
$ git merge feature/speed
Merge made by the 'ort' strategy.
 walk.py | 2 +-
 1 file changed, 1 insertion(+), 1 deletion(-)
$ git log --oneline --graph
*   4d84182 Merge feature/speed, use 0.75
|\
| * 8cc1a09 Increase speed factor
* | 3fa9bd2 Add logging to walk.py
|/
* 1e0a5c6 Add walking step function
\end{verbatim}
\begin{center}
\begin{tikzpicture}[c/.style={draw,circle,thick,minimum size=0.9cm,fill=green!12,font=\scriptsize\ttfamily},lb/.style={draw,rounded corners,font=\footnotesize\ttfamily,fill=yellow!30},>={Stealth[length=2.2mm]}]
  \begin{scope}
    \node[font=\footnotesize\bfseries] at (2.2,2.0) {fast-forward};
    \node[c] (a) at (0,0.6) {C}; \node[c,fill=orange!25] (b) at (1.5,0.6) {D}; \node[c,fill=orange!25] (cc) at (3,0.6) {E};
    \draw[->,thick] (b)--(a); \draw[->,thick] (cc)--(b);
    \node[lb] (m) at (3,-0.6) {main (now)}; \draw[->,thick] (m)--(cc);
    \node[lb,fill=gray!15] at (0,-0.6) {main (before)};
  \end{scope}
  \begin{scope}[xshift=7.2cm]
    \node[font=\footnotesize\bfseries] at (2.4,2.0) {merge commit};
    \node[c] (a) at (0,0.6) {C}; \node[c,fill=orange!25] (d) at (1.5,1.2) {D}; \node[c] (f) at (1.5,0) {F};
    \node[c,fill=blue!15] (mm) at (3.3,0.6) {M};
    \draw[->,thick] (d)--(a); \draw[->,thick] (f)--(a); \draw[->,thick] (mm)--(d); \draw[->,thick] (mm)--(f);
    \node[lb] (m) at (3.3,-0.6) {main}; \draw[->,thick] (m)--(mm);
  \end{scope}
\end{tikzpicture}
\captionof{figure}{Left: fast-forward, the label simply moves. Right: a merge commit \texttt{M} with two parents joins two lines of work.}\label{fig:g2_merge}
\end{center}

\section*{Conflicts}
If both branches changed the \emph{same lines}, Git cannot decide which version is right. This is a \textbf{merge conflict}. It is a normal event, not an error. Git stops and marks the problem:
\begin{verbatim}
$ git merge feature/speed
Auto-merging walk.py
CONFLICT (content): Merge conflict in walk.py
Automatic merge failed; fix conflicts and then commit the result.
$ git status -s
UU walk.py
\end{verbatim}
Open the file. Git wrote both versions into it:
\begin{verbatim}
def step(speed):
<<<<<<< HEAD
    return speed * 0.5   # main
=======
    return speed * 0.75  # feature/speed
>>>>>>> feature/speed
\end{verbatim}
Everything between \verb!<<<<<<< HEAD! and \verb!=======! is the version of your current branch, the rest down to \verb!>>>>>>>! comes from the other branch. Resolve it in four steps:
\begin{enumerate}[nosep]
  \item Decide what the correct code is (one version, the other, or a mixture; talk to the author if unsure).
  \item Edit the file so that it contains exactly that, and delete the three marker lines.
  \item \verb!git add walk.py! (this tells Git ``resolved'').
  \item \verb!git commit! (the default message is fine).
\end{enumerate}
If you are lost, \verb!git merge --abort! returns everything to the state before the merge. VS Code shows the two sides with buttons ``Accept Current / Incoming / Both''. Run your tests after resolving: the file may look right and still not work.

\begin{keyidea}[How to have fewer conflicts]
Conflicts grow with the time that two branches live apart and with the amount of code that both touch. So keep branches short (days, not months), merge \texttt{main} into your branch often, make small commits, and do not reformat files you do not work on.
\end{keyidea}

\section*{Rebase: a straight history}
There is a second way to combine work. \texttt{git rebase main} (on your feature branch) takes your commits, and \emph{replays} them on top of the newest \texttt{main} (Figure~\ref{fig:g2_rebase}). The result is a straight line, as if you had started your work today.
\begin{verbatim}
$ git switch feature/log
$ git rebase main
Successfully rebased and updated refs/heads/feature/log.
$ git switch main
$ git merge feature/log            # now a fast-forward
\end{verbatim}
\begin{center}
\begin{tikzpicture}[c/.style={draw,circle,thick,minimum size=0.9cm,fill=green!12,font=\scriptsize\ttfamily},>={Stealth[length=2.2mm]}]
  \node[c] (a) at (0,0.6) {C}; \node[c,fill=orange!25] (d) at (1.5,1.2) {D}; \node[c] (f) at (1.5,0) {F};
  \draw[->,thick] (d)--(a); \draw[->,thick] (f)--(a);
  \draw[->,very thick,gray!60!black] (3.0,0.6) -- node[above,font=\footnotesize\ttfamily]{rebase} (4.6,0.6);
  \node[c] (a2) at (5.5,0.6) {C}; \node[c] (f2) at (7,0.6) {F}; \node[c,fill=orange!25,dashed] (d2) at (8.5,0.6) {D$'$};
  \draw[->,thick] (f2)--(a2); \draw[->,thick] (d2)--(f2);
\end{tikzpicture}
\captionof{figure}{Rebase: commit D is copied on top of F. The copy D$'$ has a \emph{new hash}, because its parent is different.}\label{fig:g2_rebase}
\end{center}
\begin{keyidea}[The golden rule of rebase]
Rebase only commits that nobody else has yet: your own local, unpublished work. A rebase creates \emph{new} commits. If others already based work on the old ones, their history and yours no longer fit together, and a mess follows. \textbf{Never rebase a branch that you have pushed and that others use, and never rebase \texttt{main}.}
\end{keyidea}
Merge preserves the history as it happened; rebase makes it easy to read. Both are fine; teams choose one (Chapter~\ref{ch:git4}). Conflicts during a rebase are solved the same way, per commit: fix the file, \verb!git add!, \verb!git rebase --continue! (or \verb!git rebase --abort!).

\subsection*{Cleaning up your own commits}
\verb!git rebase -i HEAD~3! opens an editor with your last three commits. By changing the word in front of a line you can \texttt{reword} a message, \texttt{squash} a commit into the previous one, \texttt{drop} it, or reorder lines. This is how you turn ``wip'', ``fix typo'', ``fix again'' into one clean commit before others see it.

\section*{Stash: put work aside}
You are in the middle of an edit when you must urgently look at another branch. Your changes are not ready for a commit:
\begin{verbatim}
$ git stash                        # save changes, clean working dir
Saved working directory and index state WIP on main: dc7e3fc ...
$ git switch fix/camera            # do the urgent work
$ git switch main
$ git stash list
stash@{0}: WIP on main: dc7e3fc Ada: update README
$ git stash pop                    # bring the changes back
\end{verbatim}
\texttt{git stash -u} also stores untracked files. The stash is local and easy to forget: do not use it as long-term storage; a work-in-progress commit on a branch is safer.

\section*{Cherry-pick, tags and comparing}
\begin{description}
  \item[Cherry-pick:] copy \emph{one} commit to the current branch, for example a bug fix that is also needed on the release branch: \verb!git cherry-pick 8cc1a09!.
  \item[Tags:] a permanent name for a commit, for releases. \verb!git tag -a v0.1 -m "First walking release"! creates an annotated tag. Show with \texttt{git tag}, publish with \verb!git push origin v0.1! (Chapter~\ref{ch:git3}).
  \item[Compare:] \verb!git diff main..feature/speed! shows what the feature adds; \verb!git log main..feature/speed --oneline! lists the commits that are in the feature but not in \texttt{main}; \verb!git diff --stat main! gives a summary.
\end{description}

\begin{example}
Ada works on \texttt{feature/log} and Bob on \texttt{feature/speed}, both change line 2 of \texttt{walk.py}. Bob merges first. When Ada runs \texttt{git merge main} on her branch she gets a conflict in line 2. She opens the file, sees her \texttt{return speed * 0.5 \# log} between \texttt{HEAD} and \texttt{=======} and Bob's \texttt{0.75} below. She keeps Bob's value and her log text, deletes the markers, runs \texttt{python walk.py}, then \verb!git add walk.py! and \verb!git commit!.
\end{example}

\section*{Practice problems}
\begin{enumerate}[label=\textbf{\thechapter.\arabic*}]
  \item Which command creates the branch \texttt{fix/camera} and switches to it? Which command lists all branches?
  \item What is a branch, technically? What happens to the labels when you make a commit on the current branch?
  \item Explain the difference between a fast-forward merge and a merge commit. When can a fast-forward not happen?
  \item You run \texttt{git merge feature} and see \texttt{CONFLICT (content)}. Write the four steps to finish the merge. How do you cancel instead?
  \item In a conflicted file, which part is yours and which is the other branch's?
  \item Why does a rebased commit get a different hash?
  \item State the golden rule of rebase and give a reason for it.
  \item You have uncommitted edits and must switch branches. Which commands save and restore them?
  \item Write the command that lists the commits that are on \texttt{feature/gait} but not on \texttt{main}.
  \item[\textbf{\thechapter.10}$\star$] A bug fix is on \texttt{main} as commit \texttt{9f3c2aa}. The branch \texttt{release/1.0} needs the same fix, but not the other new work of \texttt{main}. What do you do?
\end{enumerate}
